% =====================================================================
%  §8  Errata
% =====================================================================
\section{The day's errata: 26 normalization, coordinate, convention and enclosure errors}
\label{sec:errata}

The campaign caught 26 instances of ``a quantity placed where it does not belong'' in a single day.
They are not a collection of typos but twenty-six instances of one family: \textbf{putting a quantity
where it does not belong} (the lesson of \code{LOOP-RUNBOOK.md} \S10.17, quoted verbatim). Several of
them decided whether a verdict could be pronounced at all, so each is written as
``wrong $\to$ right $\to$ source'', with a marker for the ones that changed a verdict.

\vspace{2pt}
{\footnotesize
\setlength{\tabcolsep}{3pt}
\renewcommand{\arraystretch}{1.05}
\begin{longtable}{@{}>{\raggedleft\arraybackslash}p{0.45cm} >{\raggedright\arraybackslash}p{1.95cm} >{\raggedright\arraybackslash}p{4.15cm} >{\raggedright\arraybackslash}p{4.9cm} >{\raggedright\arraybackslash}p{3.65cm}@{}}
\toprule
\textbf{\#} & \textbf{Kind} & \textbf{Wrong} & \textbf{Right} & \textbf{Source}\\
\midrule
\endfirsthead
\toprule
\textbf{\#} & \textbf{Kind} & \textbf{Wrong} & \textbf{Right} & \textbf{Source}\\
\midrule
\endhead
\bottomrule
\endfoot
1 & normalization & the number $c=12$--$18$ for $m/\Lambda$ & $c=6.3(3)$; main cause: the corpus value
$ma=1.073(88)$ was taken as the physical $0^{++}$ (too large by $1.89\times$); secondary cause:
$\Lambda$ taken as $200$--$240$ MeV instead of $263(4)(9)$ & \code{L2-GLUEBALL-LAMBDA.md};
\code{LOOP-RUNBOOK.md} \S10.16 \textbf{(verdict changed)}\\
2 & convention & $m/\Lambda=21$--$25$ (\code{a10}) & withdrawn: two independent errors $\times1.889$ and
$\times0.558$, net $3.38$ (matching the measured ratio) & \code{LOOP-RUNBOOK.md} \S10.17
\textbf{(verdict changed)}\\
3 & domain & using two-loop running to invert $\Lambda_{\mathrm{MS}}$ at $g^{2}=1$
($\beta=6$, $\mathrm{SU}(3)$) & invalid: gives $147$ MeV against the physical $263$ MeV (off by
$1.79\times$); the two-loop formula is right, the domain is wrong & \code{LOOP-RUNBOOK.md} \S10.17\\
4 & coordinate & $\lambda\sim1$ (Balaban's stopping point, reading $g^{2}$) and $\tH=1$ (lower edge of the
crossover, reading $\beta N$) treated as the same $1$ & same name, different coordinates, differing by
exactly $2N^{2}$ ($\mathrm{SU}(3)$: $18$; $\mathrm{SU}(2)$: $8$); one must cross $2.47\times$
(physical point) / $18\times$ (lower edge), not $7.3\times$ & \code{BC-CHAIN.md}:90;
\code{LOOP-RUNBOOK.md} \S9.8\\
5 & numerical & $\Gpp(\beta=2.5)=+0.5427$ & $+0.135995$: the second-order term of $\sin(\rho/2)u$, namely
$s\cos(r/2)/(2r)$, was double counted; $u''=0$ and $\rho'=0$ make it vanish. The test is the archived
$E[H]\equiv0$ & \code{H1A-EXACT.md} \S2.4\\
6 & normalization & $Q(\infty)=0.1099$ & $0.2667$ (since $0.1099=0.32957/\lambda_1$, a double error from
dividing the 13-point trapezoid by $\lambda_1$); plus $Q\propto\beta^{1.14}$, $p=0$, and the threshold
$e^{-Q}\approx0.77$ instead of $0.90$ & \code{A3-Q-SCALING.md}; \code{LOOP-RUNBOOK.md} \S10.8
\textbf{(verdict changed)}\\
7 & norm & reading $L^{4}$ as ``curvature blow-up'' & $L^{4}$ is the matrix dimension: bounding the
operator norm by the trace loses exactly $\dim=3E=12L^{4}$; measured $\mathrm{tr}/$row-sum
$=126/1642/10364$ with row sums $=3.51/4.31/3.40=O(1)$ & \code{DENSITY-D.md}; \code{LOOP-RUNBOOK.md} \S9.4\\
8 & sup/inf & using $d=9.375$ as a $\sup$ & it is a \emph{lower} bound: the correct interval is
$d\in[3.62,6]\beta^{2}$ ($\beta=2.5$: per plaquette $[22.64,37.5]$); exponents and verdicts unchanged,
the constant was underestimated by $\ge2.2\times$ & \code{R3-L4.md}; \code{LOOP-RUNBOOK.md} \S9.9\\
9 & normalization & $\operatorname{Hess}\succeq-(4.8/N)\beta$; later $-(2.9/N)\beta$ & the first is the
graph-cardinality bound without symmetrization ($\sup\lambda_{\max}=0.752$--$0.763$); the second
symmetrizes ($\tfrac12[\lambda,\lambda]$) but still sums per plaquette (missing cross-link blocks, a
factor $8/3$) $\Rightarrow$ the truth is $-8\beta/N$; consequently $\dot\lambda_0$ changes from $-3.51$
to $-6.29$, threshold $\beta<N^{2}/64$ & \code{AT-STAR.md} E3; \code{MC-METRIC.md} E1;
\code{LOOP-RUNBOOK.md} \S9.7\\
10 & coordinate & $120\times$ and $13.3\times$ (distance to the SZZ radius) & $\beta_W=N^{2}\beta_A$;
correct values $27.6\times(\mathrm{SU}(2))$ / $30.4\times(\mathrm{SU}(3))$, i.e.\ about $30\times$.
Consequently $1/32\to8.9\times$ should read $19$--$20\times$, and $\tH<N/48$, $\lambda\approx7.3$ should
read $\lambda>96$, $\lambda\approx3.2$ & \code{BETA-CONVENTION.md}; \code{LOOP-RUNBOOK.md} \S9.8
\textbf{(verdict changed)}\\
11 & normalization & using $u_{P3}=c_1/c_0$ as the per-tile weight & $u_{\mathrm{tile}}=I_2/I_1=u_{P3}/2$;
hence $m_1a(2.5)=+2.715438$, not $-0.057151$ & \code{R4-PERCENT.md} E7; \code{LOOP-RUNBOOK.md} \S9.3\\
12 & convention & $\beta_{\mathrm{tube}}=2.446918$ and ``off by $2.12\%$'' & withdrawn: a threefold
misalignment ($u_{P3}$ normalization $\times$ $|X|=1$ $\times$ threshold $1\neq a$); the direction is not
``$2\%$ off'' but \textbf{$111\%$ over}; correct $\beta_c^{KP}\le1.182227$ (per site), optimal $a=1/6$.
Also the tabulated $\varphi(\text{cube})$ $2.03\times10^{-3}$ should be $1.916\times10^{-3}$ &
\code{R4-PERCENT.md}; \code{LOOP-RUNBOOK.md} \S9.3; \code{AUDIT-R5B.md} \textbf{(verdict changed)}\\
13 & combinatorics & plaquette adjacency number $\Delta=12$ (a 3D value) & $\Delta=20$ ($4(2d-3)$,
verified by brute force for $d=2,3,4,5\to4,12,20,28$); the $\kappa$ upper bound moves from $48$ to $76$,
so $\kappa\le4.3307$ and the interval is $[1.4012,4.3307]$; also $\inf_A a_A^{1/A}$ was measured to be
increasing, so it should read $\lim$ & \code{Q5-KAPPA.md}; \code{LOOP-RUNBOOK.md} \S11.40
\textbf{(verdict changed)}\\
14 & dimension & $\mu^{2}=\delta^{3}_{\perp}(0)$ ($L^{-3}$) read directly as $\mu^{2}$ ($L^{-2}$) &
$\mu^{2}=\varepsilon\,\delta^{3}_{\perp}(0)=1/\varepsilon^{2}$; exactly one power of $a$ is missing (the
curve delta normalization $\delta_C=a^{-1}$), and it is not a convention that may be dropped &
\code{MU-3-RENORM.md} \S1; \code{MU-4-LATTICE.md} \S1.3, \S4\\
15 & sign & \code{(8.4)} giving $G_R=4\pi g^{2}NR\,G/\varepsilon^{2}$ coexisting with
$\sigma=+2g^{2}N/\varepsilon^{2}$ from \code{(8.5)} & they differ by a sign: either \code{(8.4)} needs a
minus on the right, or $(\star)$ does; the sign of $(\star)$ is not fixed & \code{MU-3-RENORM.md} \S1.2\\
16 & dimension & the one-dimensional reduction passing $f''$ off as $\partial_\mu\delta/\delta\sigma$
($L^{-2}$) & the operator is $L^{-3}$ as a whole; the power-law exponent moves from $p=-1$ to $p=0$
(which voids the ``rigid exact solution'' of $c_3$) & \code{SP-A-SPDE.md} \S0; \code{RV-3-COVARIANCE.md}
\S2; \code{LOOP-RUNBOOK.md} \S11.3\\
17 & convention & $0.2688$ taken as $\Gpp$ & it is $1/\operatorname{Var}(a)$ ($\operatorname{Var}=3.719824$);
$\Gpp$ at $\beta=0$ is $1/(2\,\mathrm{Cin}\,2\pi)=0.20511530$. The multiplicative reading
$D/S:=1-\Gpp\operatorname{Var}(a)$ gives $0.2370$ at $\beta=0$ while the true defect ratio is $0$ &
\code{DS-RATIO.md} (line 209)\\
18 & dimension & $(\star_1')$ right side $c_1f''+c_3f'/R=c_2g^{2}Nf^{2}/R^{3}$ & it should be
$c_2g^{2}Nf^{2}/R^{2}$ (one $R$ too many); this error alone voids the ``rigid exact solution'', whose
amplitude is an artefact of the wrong scaling & \code{C3-COEFF.md} \S1\\
19 & geometry & taking the intersection angle of the figure-eight to be $\pi/2$ & coplanarity plus a
single common point forces tangency, with oriented angle $\pi$; $b_0=2\cos(\varphi/2)=0$ $\Rightarrow$
RHS $\equiv0$, and $\rho_2=0$ (the $\partial^{2}/\partial R_2^{2}$ slot is permanently empty) &
\code{FIG8.md} \S1.1, \S1.4, \S2.3 \textbf{(verdict changed)}\\
20 & normalization & \code{na2}'s chain averages ``decreasing in $L$'' & $X_{na2}=X_{\text{true}}/L^{2}$
(too small by $L^{2}$); after rescaling $|X|\propto L$ and $X/L\to-0.010$ ($\beta=0.5$); the conclusion
that all $15/15$ are negative is unchanged & \code{CT-CHAIN.md} \S2.1; \code{NA2-NONDEGENERATE.md}\\
21 & citation & ``Osterwalder--Seiler 1978 proved area law $\Rightarrow$ gap'' & the abstract proves only
Wilson's confinement bound and a positive self-adjoint transfer matrix, and \textbf{contains no mass gap};
the proved implication runs $m>0\Rightarrow$ confinement (Chatterjee 2021) & \code{MU5-SIGMA-M.md}
\S1.1--\S1.2 \textbf{(verdict changed)}\\
22 & normalization & $\tH<1/24$ (as the SZZ-provable region) & no such display exists in the source; the
correct reading is $\tH<N/48$ (i.e.\ $\lambda=g^{2}N>96$), and the old $\lambda\approx7.3$ is wrong &
\code{LIT-CHECK.md} Q3$'$; \code{AUDIT-R1.md} \S A2; \code{LEDGER.md} \#7\\
23 & bookkeeping & $\sigma a^{2}=\ln(1/u)$; $\sigma>\kappa/a^{2}$ & the $-\kappa$ is missing: correctly
$c(\beta)=-\log u=\kappa+\sigma a^{2}/\pi$; the version with $\pi$ is wrong & \code{Q4-LOGCONVEX.md};
\code{LOOP-RUNBOOK.md} \S11.38(4)\\
24 & normalization & $a^{-1}(\beta=6)=2.91$ GeV (used as the anchor of \code{mu-3}/\code{mu-4}) & $2.243$
GeV (AT2021 Tab.13); high by $1.30\times$; the direction of the conclusion is unchanged but every number
that used $2.91$ must be divided by $1.30$ & \code{MU6-NUM.md}; \code{LOOP-RUNBOOK.md} \S11.37\\
25 & enclosure & $\mathrm{gap}(0)=0.729924127576750\ldots$ with an interval of width $1.4\times10^{-10}$,
quoted as a rigorous enclosure & $0.729924128620548449$ with width $\pm10^{-55}$: the four molecular
integrals and their Cauchy tails are valid enclosures, but the quotients $E[g]$, $E[\Phi'']$ were printed
with collapsed endpoints and the tails never propagated through the division, so the interval did
\emph{not} contain the truth (it lay $9.05\times10^{-10}$ outside, $6.5$ times its own width); the same
defect exceeded the $\theta=\pi/2$ width by $50\times$. Found by an independently written interval class,
and flagged independently from a second pair of coordinates & \code{R10-AUDITCERT.md} \S0, \S3;
\code{R06-GAPMONO.md} \S4 \textbf{(enclosure corrected)}\\
26 & bookkeeping & $L=2$ described as a $1\times1\times1\times1$ torus with $16$ links and $2$
plaquettes & $4L^{4}=64$ links, $6L^{4}=96$ plaquettes, $16$ sites; the $16$-link object is the blocked
lattice $L_{c}=1$ obtained from $L=2$ by a factor-$2$ blocking. The dimension count $192-48-1=143$
follows the corrected bookkeeping & \code{R04-L2CERT.md} \S1.0, with \code{O24-L2-PUSH.md}:49 as the
independent confirmation; the text corrected is \code{Q11-CERT.md} \S6\\
\end{longtable}
}

\subsection*{How to use this table (three rules)}

\begin{enumerate}[label=(\roman*)]
\item The seven rows marked ``(verdict changed)'' (1, 2, 6, 10, 12, 19, 21) are the ones that really
      affect conclusions; row 25 changes no sign but withdraws an enclosure, which for a certificate is
      the same kind of damage. The rest affect only numbers or wording. Every row was checked line by
      line, and none of them changes the sign verdicts of \S\ref{sec:tension} or \S\ref{sec:loop}. \LT
\item \textbf{Rows 4 and 10 are two attacks of the same disease}: there are two $\beta$ coordinates
      ($\beta_W=N^{2}\beta_A$). Wherever a bare $1/48$, $1/24$ or $1/32$ appears, first ask ``which
      $\beta$ is this?''. \LT
\item \textbf{Rows 14, 16 and 18 are another form of the same disease}: one power of $L$ is dropped or
      added. Wherever $\delta$, $\delta_C$, $\mu^{2}$ or $\partial_\mu\delta/\delta\sigma$ appears,
      write the slots down before computing. \LT
\end{enumerate}

\subsection*{One conclusion-level correction, made the same day}

The table above collects normalization, coordinate and convention errors only. One \emph{conclusion}
was also withdrawn on the same day and is recorded here rather than in the table: the inference
``$\beta_c(\mathrm{SU}(3))>\beta^{*}$ implies H1 fails before the $\mathrm{SU}(3)$ scaling window opens,
hence H1 is the wrong target for the Clay object'' assumed $\beta^{*}$ to be $N$-independent. The $N=3$
measurement of \eqref{eq:n3} shows $V_0(3)<0$, so $\beta^{*}(3)$ does not exist and H1 holds for every
$\beta>0$ at $N=3$ --- in particular on the whole $\mathrm{SU}(3)$ scaling region. The structural law of
\eqref{eq:n32} puts the zero at $\beta^{*}(3)\approx19.6$, i.e.\ $\beta^{*}/\beta_c=1.87\to3.45$; the
monotonicity of the measured $\delta$ only pushes it further out. \LE
\src{\code{LOOP-RUNBOOK.md} \S11.39 (lines 1879--1885) for the withdrawn conditional; \eqref{eq:n3} and
\S\ref{sec:failure} for its falsified antecedent. The measurement is a single-plaquette/patch Monte
Carlo, not the full theory.}

\vspace{3pt}
\noindent
\textbf{One lesson, quoted:} all twenty-six are the same error, ``putting a quantity where it does not
belong''. The operational form used from now on is: \emph{state the slot, then fill in the numbers}. \LT
\src{The lesson is verbatim from \code{LOOP-RUNBOOK.md} \S10.17; the ledger is \code{LEDGER.md} (50 main
rows, 10 contradictions, 11 quantities of undetermined dimension). This table collects only rows whose
``wrong $\to$ right'' can be checked character by character; the contradictions and the
undetermined-dimension items of \code{LEDGER.md} are not counted among the 26. Row 25 is the only row of
the 26 that concerns an \emph{enclosure} rather than a number, and it is recorded here rather than in the
certificate's own section because the failure mode --- a bound quoted wider than it was --- is the same
failure mode as the other twenty-five.}

